\subsection{Архитектура библиотеки и стек технологий}

\subsubsection{Общая архитектура библиотеки}

Разрабатываемая библиотека организована по модульному принципу:
каждый компонент несёт единственную ответственность и может быть
заменён или расширен независимо от остальных. Библиотека состоит
из трёх основных модулей, взаимосвязь которых отражена в структуре
заголовочных файлов.

Первый модуль --- \texttt{WorkStealingDeque} --- реализует
двустороннюю очередь задач, являющуюся ключевой структурой данных
механизма work-stealing. Данный модуль не зависит ни от каких
других компонентов библиотеки и может использоваться независимо.

Второй модуль --- \texttt{SimpleThreadPool} --- реализует
классический пул потоков с единой глобальной очередью задач.
Данный модуль использует стандартные примитивы синхронизации
C++11: \texttt{std::mutex}, \texttt{std::condition\_variable}
и \texttt{std::atomic}.

Третий модуль --- \texttt{WorkStealingPool} --- реализует пул
потоков с механизмом work-stealing. Данный модуль использует
\texttt{WorkStealingDeque} как внутреннюю структуру данных ---
по одной очереди на каждый рабочий поток.

Оба класса пулов потоков предоставляют идентичный внешний
интерфейс: метод \texttt{submit()} для постановки задач и метод
\texttt{thread\_count()} для получения числа рабочих потоков.
Идентичность интерфейса обеспечивает взаимозаменяемость реализаций
и позволяет проводить корректное сравнительное исследование
в одинаковых условиях.

Физическая структура проекта организована следующим образом.
Каталог \texttt{include/thread\_pool/} содержит заголовочные
файлы всех трёх модулей. Каталог \texttt{src/} содержит файлы
реализации для обоих пулов потоков. Каталог \texttt{tests/}
содержит наборы автоматических тестов корректности. Каталог
\texttt{benchmarks/} содержит тесты производительности.
Такое разделение соответствует общепринятым соглашениям
о структуре проектов на языке C++~\cite{Williams2019}.

\subsubsection{Обоснование выбора стека технологий}

Выбор инструментов разработки определялся тремя критериями:
соответствием стандарту C++17, отсутствием внешних зависимостей
в основной библиотеке и наличием развитой экосистемы для
тестирования и измерения производительности.

Язык C++ версий C++11/14/17 является основным инструментом
разработки. Стандарт C++11 ввёл полноценную модель памяти
для многопоточных программ и набор примитивов: \texttt{std::thread},
\texttt{std::mutex}, \texttt{std::condition\_variable},
\texttt{std::atomic}, \texttt{std::future}~\cite{Williams2019}.
Стандарт C++17 добавил \texttt{std::invoke\_result} и
\texttt{std::apply}, использованные в реализации метода
\texttt{submit()}. Все необходимые примитивы доступны в рамках
стандартной библиотеки без внешних зависимостей, что обеспечивает
максимальную переносимость библиотеки.

Система сборки CMake версии 3.14 и выше выбрана как
кроссплатформенный стандарт де-факто для проектов на C++.
CMake обеспечивает воспроизводимую сборку на различных
платформах и компиляторах, а также предоставляет механизм
\texttt{FetchContent} для автоматического подключения сторонних
зависимостей --- библиотек тестирования и измерения
производительности.

Библиотека Google Test версии 1.14.0 используется для
автоматического тестирования корректности. Выбор обусловлен
широкой распространённостью в проектах с открытым исходным
кодом, богатым набором макросов для проверки утверждений
и встроенной поддержкой в системе сборки CMake через
\texttt{gtest\_discover\_tests()}.

Библиотека Google Benchmark версии 1.8.3 используется для
измерения производительности. Она обеспечивает статистически
корректные измерения: автоматически определяет необходимое
число итераций для достижения заданной точности, вычисляет
среднее время, стандартное отклонение и позволяет выводить
результаты в формате JSON для последующего анализа.
Существенным преимуществом является поддержка параметризованных
бенчмарков через метод \texttt{Range()}, что позволяет
автоматически запускать один и тот же сценарий нагрузки
с различным числом потоков.

Сравнительная характеристика выбранных инструментов
представлена в таблице~\ref{tab:stack}.

\begin{table}[h]
\caption{Стек технологий разрабатываемой библиотеки}
\label{tab:stack}
\begin{tabularx}{\textwidth}{|l|l|X|}
\hline
\textbf{Инструмент} & \textbf{Версия} & \textbf{Назначение} \\
\hline
C++ & 11/14/17 & Язык реализации, модель памяти,
                 примитивы многопоточности \\
\hline
CMake & $\geq$ 3.14 & Система сборки, управление
                      зависимостями \\
\hline
Google Test & 1.14.0 & Автоматическое тестирование
                       корректности \\
\hline
Google Benchmark & 1.8.3 & Измерение и сравнение
                           производительности \\
\hline
\end{tabularx}
\end{table}

\subsubsection{Проектирование программного интерфейса библиотеки}

Программный интерфейс обоих пулов потоков спроектирован по
единому образцу. Центральным элементом интерфейса является
шаблонный метод \texttt{submit()}, принимающий произвольный
вызываемый объект и возвращающий \texttt{std::future}
с результатом его выполнения:

\begin{lstlisting}[language=C++, caption={Интерфейс метода submit}]
template<typename F, typename... Args>
auto submit(F&& f, Args&&... args)
    -> std::future<typename std::invoke_result<F, Args...>::type>;
\end{lstlisting}

Шаблонные параметры \texttt{F} и \texttt{Args...} позволяют
передавать в пул любые вызываемые объекты: свободные функции,
лямбда-выражения, объекты с перегруженным \texttt{operator()}.
Тип возвращаемого значения автоматически выводится посредством
\texttt{std::invoke\_result}, что избавляет пользователя от
необходимости явно указывать его.

Возвращаемый объект \texttt{std::future} предоставляет три
возможности: получить результат выполнения задачи методом
\texttt{get()} (с блокировкой до завершения), проверить
готовность результата без блокировки методом \texttt{wait\_for()}
и получить исключение, если оно возникло в ходе выполнения
задачи.

Управление жизненным циклом пула осуществляется через конструктор
и деструктор. Конструктор принимает опциональный параметр ---
число рабочих потоков. Если параметр не указан, используется
значение \texttt{std::thread::hardware\_concurrency()},
соответствующее числу логических ядер процессора. Деструктор
обеспечивает корректное завершение: он устанавливает флаг
остановки, пробуждает все ожидающие потоки и дожидается их
завершения посредством \texttt{std::thread::join()}.

Оба класса объявлены некопируемыми и неперемещаемыми
посредством явного удаления соответствующих специальных
функций-членов. Данное решение обусловлено тем, что объект
пула потоков управляет ресурсами операционной системы
(потоками и мьютексами), семантика копирования и перемещения
которых неоднозначна.